\documentclass[12pt]{article}

\usepackage[utf8]{inputenc}
\usepackage{amsmath, amssymb, amsthm}
\usepackage{geometry}
\geometry{a4paper, margin=2.5cm}

\title{Funcții Olomorfe}
\author{}
\date{}

\theoremstyle{definition}
\newtheorem{definition}{Definiție}

\theoremstyle{plain}
\newtheorem{proposition}{Propoziție}

\begin{document}

\maketitle

\begin{definition}
Fie $\Omega \subset \mathbb{C}$ o mulțime deschisă și $f:\Omega \to \mathbb{C}$. 
Spunem că $f$ este \emph{derivabilă complex} în punctul $z_0 \in \Omega$ dacă există limita


\[
f'(z_0) = \lim_{h \to 0} \frac{f(z_0 + h) - f(z_0)}{h}.
\]


Funcția $f$ se numește \emph{olomorfă} în $\Omega$ dacă este derivabilă în fiecare punct al lui $\Omega$.
\end{definition}

\begin{proposition}
Fie $f:\Omega \to \mathbb{C}$ și $z_0 \in \Omega$. Următoarele afirmații sunt echivalente:
\begin{enumerate}
    \item[(i)] Funcția $f$ este derivabilă în $z_0$, iar derivata este $f'(z_0)=a$.
    \item[(ii)] Există o funcție $\psi$, definită în vecinătatea lui $0$, cu proprietatea
    

\[
    \lim_{h \to 0} \psi(h) = 0,
    \]


    astfel încât pentru $h$ suficient de mic are loc relația
    

\[
    f(z_0 + h) - f(z_0) - a h = h \psi(h).
    \]


\end{enumerate}
\end{proposition}

\end{document}
